% ECONOMICS_PROBLEM: {"id": "C73-19", "title": "Polynomial algorithms for games with ergodic payoff", "jel_code": "C73", "source_id": "OP-0164"}
% !TeX program = xelatex
\documentclass[11pt,a4paper]{article}
\usepackage[margin=23mm]{geometry}
\usepackage{fontspec}
\setmainfont{Latin Modern Roman}
\usepackage{amsmath,amssymb,mathtools}
\usepackage{xcolor,enumitem,booktabs,longtable,array,xurl,hyperref}
\definecolor{muted}{HTML}{53636C}
\hypersetup{colorlinks=true,urlcolor=blue,linkcolor=blue}
\setlength{\parindent}{0pt}
\setlength{\parskip}{5pt}
\newcommand{\R}{\mathbb R}
\newcommand{\E}{\mathbb E}
\newcommand{\N}{\mathbb N}
\newcommand{\ind}{\mathbf 1}
\newcommand{\Var}{\operatorname{Var}}
\newcommand{\Cov}{\operatorname{Cov}}
\newcommand{\supp}{\operatorname{supp}}
\DeclareMathOperator*{\argmin}{arg\,min}
\DeclareMathOperator*{\argmax}{arg\,max}
\newcommand{\entrymeta}[3]{{\sffamily\small\color{muted}\textbf{\texttt{#1}}\quad|\quad #2\par #3\par}}
\newcommand{\Problemref}[1]{\texttt{#1}}
\newcommand{\JELref}[2]{\texttt{#2} (JEL \texttt{#1})}
\begin{document}
\section*{Polynomial algorithms for games with ergodic payoff}
\textbf{Problem ID:} \texttt{C73-19}\quad \textbf{JEL:} \texttt{C73}\par
% BEGIN ECONOMICS STATEMENT
\label{problem:OP-0164}
% GAME_THEORY_PROBLEM: {"local_id": "GT-034", "original_number": 34, "kind": "R", "entry_kind": "statement_only", "published": false, "review_required": true, "category": "game_theory"}
\label{game:GT-034}
{\sffamily\small\color{muted}
{\sffamily\small\color{muted}\textbf{GT-034} | Stochastic games and computational complexity\par R | Representative formalization of a broad source agenda\par}
\par}
\subsection*{Definitions and mathematical statement}
One precise representative problem uses finite, turn-based, zero-sum stochastic mean-payoff games. An input explicitly lists a finite state set, the controlling player at each state, finite actions, rational transition probabilities, bounded rational rewards $r(s,a)$ to Max, and a rational threshold $q$, all in binary. For arbitrary behavioral strategies define
\[
 g_s(\sigma,\tau)=\mathbb E_{s,\sigma,\tau}
 \left[\liminf_{T\to\infty}\frac1T\sum_{t=0}^{T-1}r(s_t,a_t)\right],
\]
\[
 v_-(s)=\sup_\sigma\inf_\tau g_s(\sigma,\tau),
 \qquad v_+(s)=\inf_\tau\sup_\sigma g_s(\sigma,\tau).
\]
Impose the explicitly stated promise
\[
 \exists\rho\in\mathbb R\quad\forall s\in S,\qquad
 v_-(s)=v_+(s)=\rho.
\]
Is there an algorithm that, on every promised input, decides $\rho\geq q$ in time polynomial in the entire input bit length? This formulation treats state-independent value as its precise ergodicity promise.
\subsection*{Short English statement}
Can a clearly specified class of stochastic mean-payoff games with a common value across initial states be solved in polynomial time?
\subsection*{Variants and scope boundaries}
State-independent value, existence of an additive eigenvector for a Shapley operator, and irreducibility under every stationary policy pair are different possible ergodicity assumptions. The representative promise above does not assert their equivalence. Simultaneous-action games and computation of an exact real value need separate specifications. Expected pathwise liminf is fixed here; exchanging the expectation and liminf defines another evaluation.
\subsection*{Source relationship and status}
[S1] records the polynomial-algorithm question for stochastic games with ergodic payment as an agenda. The report does not there fix all of the input, payoff, ergodicity, and output conventions needed for one complexity theorem. The displayed problem is editorial and is not claimed to be the uniquely intended or equivalent source model.
\subsection*{References}
\begingroup\footnotesize\raggedright
\textbf{[S1]} Inria research team TROPICAL, \emph{TROPICAL Activity Report 2024} (2024). Locator: Section 3.1, printed p. 4, algorithmic paragraph concerning polynomial algorithms for stochastic games with ergodic payment. \url{https://radar.inria.fr/rapportsactivite/RA2024/tropical/TROPICAL-RA-2024.pdf}
\par\endgroup

\subsection*{Common conventions: Source mathematical conventions}

\textbf{Finite games.} For a finite set $A$, $\Delta(A)$ is its probability
simplex. Players choose independent mixed strategies in Nash equilibrium;
correlation is allowed only when explicitly part of the solution concept.
In a game with payoff functions $u_i$ and strategy profile $x$, an additive
$\varepsilon$ Nash equilibrium satisfies
\[
 u_i(x)\geq u_i(x_i',x_{-i})-\varepsilon
 \quad\text{for every player $i$ and every admissible deviation $x_i'$}.
\]
For numerical approximation, payoffs are normalized as locally specified.
An $\varepsilon$ payoff residual is distinct from distance at most
$\varepsilon$ to the set of exact equilibria. Point convergence, convergence
of empirical play, and convergence in distance to an equilibrium set are
also distinct.

\textbf{Computational representations.} Unless stated otherwise, a complexity
question takes rational data with binary encoding; polynomial time is measured
in total bit length. A fixed-accuracy polynomial algorithm is not necessarily
polynomial in $1/\varepsilon$, and neither is necessarily polynomial in
$\log(1/\varepsilon)$. Succinct games and value, demand, or sampling oracles
must declare their interfaces. Exact real solutions require an explicit
representation; an unspecified list of arbitrary real numbers is not a
finite computational output. A claimed algorithmic barrier is meaningful
only for its specified model and reductions.

\textbf{Dynamic games.} Histories, information, observation, and allowable
randomization are part of the model. A stationary strategy depends only on
the current state; a positional strategy in a deterministic graph game
chooses one action at each controlled vertex. Nash equilibrium at an initial
state and equilibrium after every history are different requirements.
An expectation of a pathwise $\liminf$ payoff, a $\liminf$ of expected averages,
a limiting expected average, and a uniform finite-horizon equilibrium are
different criteria. None is silently substituted for another.

\textbf{Allocation and mechanisms.} A complete allocation assigns every item
exactly once unless otherwise stated. Zero-valued goods, negative values,
capacity constraints, feasibility, lotteries, and copies of goods affect
fairness definitions and are specified locally. Dominant-strategy incentive
compatibility, Bayesian incentive compatibility, universal truthfulness,
and truthfulness in expectation impose different inequalities. Whenever
an objective ratio has zero denominator, its convention is stated or those
instances are excluded.

\textbf{Infinite-dimensional models.} Expectations, integrals, stochastic
equations, and optimization objectives require their local regularity,
integrability, filtrations, and admissible domains. A backward--forward
mean-field system is a boundary-value problem; it is not an initial-value
semiflow without an additional construction. Long-time, large-population,
and vanishing-noise limits require an order of limits and a convergence
topology. If a minimum need not be attained, the objective is its infimum
or supremum and the approximation requirement is stated separately.
% END ECONOMICS STATEMENT
\end{document}
